\section{Introduction}
\label{sec:introduction}

Repeated observations of small Solar System bodies provide positions for
orbit determination and photometric measurements of rotational variability.
These uses place different demands on an observing program: an angular
trajectory must be sampled sufficiently to recognize a moving source,
whereas physical characterization benefits from measurements repeated over
longer intervals \citep{WarnerHarrisPravec2009,Denneau2013}. For a wide-field telescope,
the combination of sky coverage, exposure time, and revisit interval
therefore determines which objects can be measured and how the observations
can be used.

Modern surveys illustrate several ways of making these choices. Pan-STARRS
couples wide-field imaging to attribution and linking across multiple
nights; ATLAS emphasizes rapid sky coverage for nearby-object searches;
and the Zwicky Transient Facility (ZTF) combines a large field of view with
repeated imaging for moving-object and time-domain science
\citep{Denneau2013,Tonry2018,Bellm2019,Masci2019}. Dedicated moving-object
processing can also exploit the cadence of a broader time-domain survey,
as demonstrated by the Heliocentric-orbiting Objects Processing System for
the Wide Field Survey Telescope \citep{Wang2025HOPS}. These programs make
clear that sky area, sensitivity, and processing must be considered
together with the temporal sampling actually delivered.

A Schmidt telescope with a moderate aperture and a large field can build a
substantial sample of asteroid measurements through repeated ecliptic
imaging. Its scientific scope is established by the observed footprint,
cadence, and measurement quality, together with the populations represented
on real images. Characterizing these quantities is necessary for selecting
observations suitable for asteroid astrometry, photometric time series,
and moving-source identification.

We developed \SHARP\ for the Xinglong 60/90~cm Schmidt telescope. Its
approximately $2.94^\circ\times2.95^\circ$ field of view is used for repeated
30~s imaging of a predefined ecliptic field library. The observing strategy
aims for three visits to a field during a night, allowing angular motion to
be sampled over roughly an hour. Calibrated source catalogs support two
complementary analyses: association with the predicted positions of known
asteroids and a search for coherent motion among catalog detections. The
initial associated population is dominated by main-belt asteroids, making
their repeated astrometry and photometry the principal initial data product.

Known objects provide a direct way to examine the measurements and
processing. Their predicted positions can be compared with observed
centroids, repeated associations reveal the available time sampling, and
their identities allow stored moving-source linkages to be evaluated on
real detections. The GOTTA Prototype analysis provides a related
application of known-asteroid measurements to survey characterization
\citep[in preparation]{Xiao2026InPrep}. We combine these tests with exposure-level depth
estimates and repeated stellar photometry. The independent motion search
adds an image-reviewed sample of single-night linkages. Linking methods
based on angular trajectories or heliocentric transformations connect such
short arcs to the longer astrometric sequences needed for identification
and orbit determination \citep{Kubica2007,Holman2018,Moeyens2021}.

This paper describes the survey design and its 2025 November--2026 July
initial observing campaign. Section~\ref{sec:survey} introduces the instrument,
photometric convention, observing strategy, and analysis samples.
Section~\ref{sec:processing} describes the moving-object analysis.
Section~\ref{sec:performance} measures the realized footprint, cadence,
photometric and astrometric behavior, and conditional linkage retention.
Section~\ref{sec:samples} presents the associated known-object population
and historical single-night candidates. We discuss the scientific scope
and subsequent use of these measurements in Section~\ref{sec:discussion}
and conclude in Section~\ref{sec:conclusions}.
